\section{Discussion and limitations}
\label{sec:discussion}
\textbf{What the results say.} Reporting the score of an automatically tuned configuration on the questions that guided the search overstates its accuracy on new questions by an amount that is large when the dev set is the size typical for tuning ($n\lesssim 200$), shrinks only slowly with $n$, and depends on the search-space diversity and the signal-to-noise ratio in the way \cref{prop:opt} predicts. The inflation spills over to the conclusions drawn from the numbers: lifts over a baseline, significance of a paired test, and comparisons between systems tuned with different budgets. Sample splitting makes the report unbiased for any optimizer at the cost of selection quality and of interval width; model-based corrections that exploit the paired structure of the outcomes (empirical Bayes, replay bootstrap) are markedly more accurate when every configuration is scored on the same questions, but their validity rests on assumptions that hold only approximately.

\textbf{Practical recommendations.} (1) Report the held-out score of the deployed configuration, with the number of configurations/trials searched and $n$; a 50/50 split is a reasonable default (\cref{tab:split}: selecting on 30\% gives a more precise report in every world but loses up to 1.3 points of deployed accuracy, selecting on 70\% a much noisier report). (2) If all questions must be used for search, report a search-aware estimate (empirical Bayes or replay bootstrap) \emph{together with} the naive score, and a simultaneous lower bound if a guarantee is needed. (3) Do not compare systems tuned with unequal budgets on a shared dev set without correcting for the budget. (4) Treat dev-set lifts smaller than the optimism predicted by \cref{prop:opt} (about $\sigma_n\sqrt{1-\rho}\,m_K$ for equally good configurations) as unresolved.

\textbf{Limitations.} \emph{Simulated search over cached configurations.} Our optimizers replay the search logic of DSPy/MIPRO/best-arm-identification methods on discrete, pre-computed spaces of 40--80 configurations; we did not run those tools end-to-end (their LLM-proposed instructions and bootstrapped demonstrations are open-ended spaces and need generation-based evaluation that was beyond our CPU budget). The qualitative mechanism does not depend on this, but absolute gaps for those tools could differ. \emph{Models and task.} The worlds use Qwen2.5-0.5B/1.5B-Instruct and SmolLM2-1.7B-Instruct (three models, two families) with log-probability answer scoring on a random sample of MMLU-Pro questions of $\le 200$ tokens (about 65\% of the benchmark; long law/history items are under-represented), without chain-of-thought and with deterministic decoding; prompt effects on larger models or generative reasoning pipelines may be larger or smaller, which by \cref{prop:opt} moves the SNR and hence the gap. \emph{Finite pool.} Truth is defined on a pool of $N{=}750$ questions per world ($N{=}450$ for world D); the benchmark is the empirical distribution of the pool, and results for $n$ close to $N$ should be read with that in mind (the without-replacement check in \cref{tab:norepl} covers $n\le 200$; bootstrapping the pool gives a 95\% range of about $\pm0.5$--$0.7$ points for the gap at $n{=}100$, \cref{tab:poolboot}). \emph{Protocols.} Empirical Bayes assumes exchangeable, roughly normal true accuracies and a fully evaluated configuration set; its simple interval loses coverage at large $n$ and for adaptive optimizers; max-$t$ and the conditional method were applied to exhaustive/random search only (applying them to adaptively visited sets is heuristic and, for max-$t$ with few visited configurations, fails); the conditional method uses a Gaussian approximation to discrete accuracies and is poorly calibrated here. We did not evaluate the reusable holdout or differentially private mechanisms beyond the Ladder, nor distribution shift between benchmark and deployment, question clustering, or label noise.
